\section{Threat Model}
\label{sec:threat-model}

\subsection{Attacker Model}
\label{sec:threat-model:attacker}

We consider a software-only adversary executing on a RISC-V Linux system.
The adversary can run arbitrary user-space code, create child processes
through \texttt{fork}, transform the address space through \texttt{execve},
and load dynamic shared objects at runtime.
The adversary may attempt anti-analysis techniques such as timing checks,
debugger detection, and process-tree manipulation to hide or manipulate
behavior and evade analysis.

\subsection{Assumptions and Trust Boundaries}
\label{sec:threat-model:assumptions}

The hardware trace adapter is part of the trusted computing base: the
adversary cannot modify, disable, or observe the trace logic.
The offline host that receives the trace buffer and performs reconstruction
is also trusted.
The target workload runs under an experiment harness that marks the start
and end of the measured region with hardware markers, defining the analysis
window.
Kernel-level rootkits that compromise the operating system itself are
outside our scope.

\subsection{Scope Statement}
\label{sec:threat-model:scope}

We focus on controlled RISC-V Linux workloads.
We do not evaluate production streaming throughput or cycle-level runtime
overhead, which we leave for future work.
